\section*{Overview}

Aliens trick is the competitive-programming name for a small Lagrangian-relaxation pattern: if the hard part of a DP is
``use at most \(K\) groups / segments / chosen objects,'' add a penalty \(\lambda\) per chosen group, solve the relaxed
problem quickly, and binary search \(\lambda\).

\section*{When to Use It}

Use it when:

\begin{itemize}
  \item the original DP has an extra dimension counting how many segments or objects were used,
  \item fixing a penalty per chosen unit collapses that dimension,
  \item the number of chosen units in the optimal relaxed solution changes monotonically with the penalty.
\end{itemize}

\section*{Core Idea}

Turn
\[
\text{maximize } \texttt{gain} \quad \text{subject to } \texttt{count} \le K
\]
into
\[
\text{maximize } \texttt{gain} - \lambda \cdot \texttt{count}.
\]

For one fixed \(\lambda\), compute both:

\begin{itemize}
  \item the best penalized value,
  \item how many groups that optimum used.
\end{itemize}

Then binary search \(\lambda\) until the chosen-count crosses \(K\), and undo the penalty at the end.

\section*{Key Insight}

The penalty parameter acts like a shadow price for opening a new group. If opening too many groups is the only thing
that makes the DP expensive, charging for each group can remove the explicit count dimension entirely.

\section*{Worked Problem}

\paragraph{Problem.}
Given an array of profits, choose at most \(K\) pairwise disjoint non-empty subarrays with maximum total sum.

\paragraph{Why the trick fits.}
The natural DP has a dimension for ``how many subarrays have been started so far.'' If starting a new subarray costs a
penalty \(\lambda\), the DP becomes linear:

\begin{itemize}
  \item either extend the current chosen segment,
  \item or start a new segment and pay \(\lambda\).
\end{itemize}

\section*{Correctness Intuition}

As \(\lambda\) increases, starting a new group becomes less attractive, so the optimal group count does not increase.
That monotonicity makes the penalty searchable. After finding the largest \(\lambda\) that still uses at least \(K\)
groups, adding back \(\lambda K\) recovers the original objective value.

\section*{Complexity Analysis}

If the relaxed solver costs \(T\), the total cost is \(O(T \log A)\), where \(A\) is the searched penalty range.

\section*{Implementation}

The code below solves the ``at most \(K\) disjoint subarrays'' problem. The relaxed solver returns a pair:

\begin{itemize}
  \item penalized score,
  \item number of used segments.
\end{itemize}

\section*{Common Pitfalls}

\begin{itemize}
  \item Applying the trick without monotonicity of the chosen count in \(\lambda\).
  \item Forgetting to define how ties compare; the segment count on equal penalized score matters.
  \item Using it when the relaxed DP is still too slow, which means the trick removed no real bottleneck.
\end{itemize}

\section*{Variants / Extensions}

\begin{itemize}
  \item Partition DPs with a penalty per segment.
  \item Resource-limited problems where one capacity is turned into a multiplier.
  \item Combining the relaxed solver with convex hull or divide-and-conquer optimization.
\end{itemize}

\section*{Practice Problems}

\begin{itemize}
  \item Maximize profit using at most \(K\) disjoint segments.
  \item Minimize cost using at most \(K\) groups after adding a per-group penalty.
  \item DP problems where one count dimension is the only reason the state explodes.
\end{itemize}

\section*{References}

\begin{itemize}
  \item \href{https://codeforces.com/catalog}{Codeforces Catalog}.
  \item \href{https://usaco.guide/adv/lagrange?lang=cpp}{USACO Guide: Lagrangian Relaxation}.
  \item \href{https://oiwiki.com/dp/opt/wqs-binary-search/}{OI Wiki: WQS Binary Search}.
\end{itemize}
